\documentclass[11pt]{article}
\usepackage[margin=1in]{geometry}
\usepackage{amsmath,amssymb,amsthm}
\usepackage[hidelinks]{hyperref}
\newtheorem{theorem}{Theorem}
\newtheorem{lemma}[theorem]{Lemma}
\newtheorem{corollary}[theorem]{Corollary}
\newcommand{\mr}{\operatorname{mr}}
\newcommand{\rank}{\operatorname{rank}}
\newcommand{\I}{\mathcal I}
\newcommand{\Q}{\mathbb Q}
\title{Additivity of normalized involutory cover bounds}
\author{Internal research working proof}
\date{17 November 2026}
\begin{document}
\maketitle

The arbitrary-cycle cover bound can lose information when it is rounded
before taking connected sums. Keeping its unrounded value gives an
exactly additive quantity for involutory meridian actions. The gluing
construction scales the number of moving pairs using a complete
bipartite graph. In particular, one of the twenty-four retained
seven-sheet sharp knots can be joined to any number of simple-cover
sharp knots without losing sharpness. This proves infinite families of
meridional-rank equalities from those certificates. These are written
working proofs, without a claim of worldwide priority, formal
verification, or external mathematical review.

\section{Normalized bounds}

For a knot $K$, consider finite transitive actions
$\rho:\pi_1(S^3\setminus K)\to S_n$, $n\ge2$, in which meridians are
nonidentity involutions. Write $Y_\rho$ for the connected branched cover.
If the meridian has $\delta$ moving pairs and $n-2\delta$ fixed sheets,
put
\[
 q(\rho)=\frac{n-1+\rank\pi_1(Y_\rho)}{\delta},\qquad
 q_k(\rho)=\frac{n-1+b_1(Y_\rho;k)}{\delta},
\]
where $k$ is a fixed field. Define the real numbers
\[
       \I(K)=\sup_\rho q(\rho),\qquad
       \I_k(K)=\sup_\rho q_k(\rho).
\]
Double covers supply nonempty sets. The retained arbitrary-cycle
bound gives
\[
       1\le \I_k(K)\le\I(K)\le\mr(K)\le b(K).
\]
Thus the suprema are finite; their attainment and rationality are not
assumed. The simple-cover rank invariant satisfies $C(K)\le\I(K)$.
The corresponding integer lower bound is $\lceil\I(K)\rceil$; it is at
most the retained arbitrary-cycle invariant $D(K)$.

\begin{theorem}\label{thm:additivity}
For any classical knots $K,L$ and any fixed field $k$,
\[
 \I(K\#L)=\I(K)+\I(L)-1,\qquad
 \I_k(K\#L)=\I_k(K)+\I_k(L)-1.
\]
The assertion concerns these normalized cover invariants. It does not
assert unrestricted additivity of meridional rank or of $D$.
\end{theorem}

\section{An arbitrary coloring of a connected sum}

Fix a transitive $n$-sheet involutory coloring of $K\#L$ and identify
the two meridians in its amalgamated knot-group presentation.
Let their common permutation have $\delta$ moving pairs and
$c=n-\delta$ cycles. The restrictions to the two summand groups may
be disconnected. Denote their sheet orbits by $O_{j\alpha}$, for
$j=1,2$, and let $N_j$ be the number of orbits. Write
$n_{j\alpha}=|O_{j\alpha}|$, $\delta_{j\alpha}$ for the number of
moving pairs in that orbit, and $Y_{j\alpha}$ for its branched cover.
If $\delta_{j\alpha}=0$, normal generation by a meridian implies
that the orbit is a singleton and its cover is $S^3$. Such an orbit
has $n_{j\alpha}-1+\rank\pi_1(Y_{j\alpha})=0$.
For each $j$,
\[
        \sum_\alpha n_{j\alpha}=n,\qquad
        \sum_\alpha\delta_{j\alpha}=\delta.
\]

The decomposing sphere intersects the knot twice. Its lift consists
of $c$ spheres: one for each fixed sheet and one double cover for each
moving pair. A trivial proper arc similarly lifts to $c$ balls.
Capping the summand covers along these balls gives the manifolds
$Y_{j\alpha}$. The graph recording their sphere gluings has
$N_1+N_2$ vertices and $c$ edges and is connected, since the original
coloring is transitive. Its first Betti number is
\[
                         \beta=c-N_1-N_2+1.
\]
Consequently
\[
 Y\cong\left(\mathop{\#}_{j,\alpha}Y_{j\alpha}\right)
          \#\left(\mathop{\#}^{\beta}S^1\times S^2\right).
\]
This is the same punctured-ball gluing argument as for the retained
simple-cover formula; it uses no geometric recognition of the
summand covers.

Van Kampen and the classical Grushko theorem \cite[Theorem 1.B.5]{stallings}
give
\[
 \rank\pi_1(Y)=\sum_{j,\alpha}\rank\pi_1(Y_{j\alpha})+\beta.
\]
Thus, writing $Q(Y)=n-1+\rank\pi_1(Y)$, one obtains the exact identity
\begin{equation}\label{eq:numerator}
 Q(Y)=\sum_{j,\alpha}
  \bigl(n_{j\alpha}-1+\rank\pi_1(Y_{j\alpha})\bigr)-\delta.
\end{equation}
For the scalar version, the same identity holds with first Betti
numbers and does not require Grushko's theorem.

For $\delta_{j\alpha}>0$, each term in the sum is at most
$\delta_{j\alpha}\I(K_j)$. Divide \eqref{eq:numerator} by $\delta$
and use the displayed sums to obtain
\[
                         q(\rho)\le\I(K)+\I(L)-1.
\]
Taking suprema proves the upper inequality in Theorem~\ref{thm:additivity}.
The scalar upper inequality follows identically with the same field
on every component.

\section{Scaling the moving pairs gives the reverse inequality}

Choose arbitrary connected involutory covers of $K$ and $L$, with
moving-pair counts $a,b$ and numerator values $Q_K,Q_L$.
Put $g=\gcd(a,b)$ and $\delta=ab/g$.
Take $b/g$ disjoint copies of the first cover and $a/g$ copies of
the second. Their meridian permutations each have $\delta$ moving
pairs. Arrange their common moving sheets as follows.

Make the complete bipartite multigraph with $b/g$ left vertices,
$a/g$ right vertices, and $g$ parallel edges between each left-right
pair. Label its $ab/g=\delta$ edges by distinct pairs of sheets.
Each left vertex has degree $a$, and each right vertex degree $b$.
Assign the moving pairs of its corresponding cover copy to its
incident edges. This makes the two meridian permutations identical.
Give every nonmoving sheet of every copy a separate additional
label, and let the other summand act trivially on that label.
Both restricted actions now live on one finite set of sheets.
They agree on the amalgamated meridian and define a coloring of
$K\#L$.

Each selected copy acts transitively on its active sheets. Two
copies at the endpoints of a multigraph edge share that moving
pair, so their active sets meet. The complete bipartite graph is
connected, including when one side has one vertex. Hence the
combined action is transitive. Its restriction components are
exactly the selected copies and the added trivial singleton covers.
By \eqref{eq:numerator},
\[
 \frac{Q(Y)}{\delta}
 =\frac{(b/g)Q_K+(a/g)Q_L-\delta}{\delta}
 =\frac{Q_K}{a}+\frac{Q_L}{b}-1.
\]
Choose the original covers arbitrarily close to the two suprema.
This gives the reverse inequality and proves the group-rank identity.
The construction and numerator identity also prove the scalar one.
No compactness, bounded-degree enumeration, or maximizing cover is
needed.

\section{A sharp connected-sum consequence}

\begin{corollary}\label{cor:onefraction}
Suppose $K$ has a transitive involutory cover whose normalized rank or
fixed-field scalar score exceeds $b(K)-1$, and let
$J_1,\ldots,J_s$ admit matching simple-cover lower and bridge upper
certificates $C(J_i)=b(J_i)=b_i$.
Then
\[
 \mr(K\#J_1\#\cdots\#J_s)
   =b(K\#J_1\#\cdots\#J_s)
   =b(K)+\sum_i(b_i-1).
\]
For the scalar version, use simple-cover scalar certificates over the
same field as the score on $K$.
\end{corollary}
\begin{proof}
Theorem~\ref{thm:additivity} and $\I(J_i)\ge C(J_i)=b_i$ give a
normalized score strictly greater than
$b(K)-1+\sum_i(b_i-1)$. Since meridional rank is an integer, its lower
bound is at least $b(K)+\sum_i(b_i-1)$. Joining bridge positions gives
the matching upper bound. The scalar assertion is identical.
\end{proof}

For each of the twenty-four retained sharp-four knots in the
seven-sheet input, the cover group is $C_3$ and its normalized
score is $7/2$. Its bridge upper bound is four. Therefore, for any
simple-cover-sharp summands as in Corollary~\ref{cor:onefraction},
\[
              \mr(K\#J_1\#\cdots\#J_s)
               =b(K\#J_1\#\cdots\#J_s)
               =4+\sum_i(b_i-1).
\]
In particular, for every such $K$ and every $s\ge0$,
\[
                      \mr(K\#^{s}T(2,3))
                       =b(K\#^{s}T(2,3))=4+s.
\]
The trefoil's double branched cover has group $C_3$, giving its
simple-cover bound two. This is an infinite equality family obtained
from the retained finite certificates, without a global priority claim.

An explicit cover avoids even taking suprema in this application.
Choose a sharp simple cover of $J_i$ on $r_i$ sheets with group $G_i$.
Use two copies on disjoint moving pairs. Align those two pairs with
the two moving pairs of the seven-sheet cover of $K$, and give all
other sheets new labels. The total degree is
\[
                         n=7+2\sum_i(r_i-2),\qquad\delta=2.
\]
The lifted sphere graph is a tree, so the cover group is
\[
               C_3*\mathop{*}_i(G_i*G_i).
\]
Its normalized score is exactly
$7/2+\sum_i(b_i-1)$, whose ceiling is the asserted equality.
For $s$ trefoils the degree stays seven and the cover group is the
free product of $2s+1$ copies of $C_3$.

For several seven-sheet sharp factors, the same invariant only gives
\[
       \mr(K_1\#\cdots\#K_t)
                    \ge\left\lceil\frac{5t+2}{2}\right\rceil.
\]
Their available bridge upper bound is $3t+1$. For $t\ge2$ the two
bounds differ. Exact normalized additivity explains this rounding
loss; it does not prove MRC for arbitrary sums of those factors.

\section{Finite controls and limits}

The accompanying checker verifies the bipartite pair assignment for
all $1\le a,b\le8$, using component intersections to test
transitivity and the sphere-graph formula to check the normalized
numerator identity. It also constructs literal presentations of the meridian-square
quotients of amalgamated knot groups from the retained diagram for
$13a_{373}$ and zero, one, two and three trefoil factors. Imposing
meridian squares first is valid because these loops already die in
the completed branched covers; fixed-sheet meridian cycles are still
filled explicitly. The full lifted-edge
homology construction gives mod-three Betti numbers $1,3,5,7$,
respectively, and resulting lower bounds $4,5,6,7$. These are finite
controls of the construction; the uniform statements require the
written gluing, supremum and classical rank arguments above.

The source knot equalities and diagram identities remain those of the
retained seven-sheet certificates. These arguments do not audit their
worldwide priority, identify all possible involutory covers, or settle
the full Meridional Rank Conjecture. In particular no assumption that
$\I(K)=7/2$ for a retained seed is made: its known cover only gives
$\I(K)\ge7/2$.

\begin{thebibliography}{9}
\bibitem{stallings} J. Stallings, \emph{Group Theory and
Three-Dimensional Manifolds}, Yale University Press (1971),
Theorem 1.B.5, printed p.~6.
\url{https://webhomes.maths.ed.ac.uk/~v1ranick/papers/stallings1.pdf}.
\end{thebibliography}
\end{document}
